\clearpage\section{Metodologia e resultados experimentais}
Foram definidos 14 ensaios no protocolo de 23/11/2026. Nos T01–T11, o Monitor Serial foi utilizado para comandos e contadores; nos T12–T14, os scripts Python transferiram imagens. A consolidação de 24/11/2026 distingue resultado relatado, comprovação parcial, sucesso funcional e integridade confirmada. As evidências dos T01–T09 presentes no relatório anterior complementam a revisão recente.

\begin{table}[H]\centering\small\caption{ID / Ensaio / Resultado / classificação}\label{tab:14}
\rowcolors{2}{tcdlight}{white}\begin{tabularx}{\linewidth}{>{\hsize=0.3093\hsize\linewidth=\hsize\raggedright\arraybackslash}X>{\hsize=1.0825\hsize\linewidth=\hsize\raggedright\arraybackslash}X>{\hsize=1.6082\hsize\linewidth=\hsize\raggedright\arraybackslash}X}
\rowcolor{tcdtablehead} \textcolor{black}{\textbf{ID}} & \textcolor{black}{\textbf{Ensaio}} & \textcolor{black}{\textbf{Resultado /\allowbreak  classificação}} \\
T01 & Autoteste CRC-8 e 4B6B & Três verificações OK; relatado. \\
T02 & Texto A \ensuremath{\to} B & Entrega e END; relatado. \\
T03 & Texto B \ensuremath{\to} A & Entrega; 2 reenvios; relatado. \\
T04 & Fragmentação & Entrega; 18 reenvios; relatado. \\
T05 & Corrupção CRC-8 & 1 CRC errado; recuperação; relatado. \\
T06 & Três erros simulados & 3 CRC errados; 1 mensagem; parcial. \\
T07 & Perda ACK DATA & Omissão, repetição e entrega; relatado. \\
T08 & Perda ACK END & END na terceira tentativa; relatado. \\
T09 & Lote curto & 10 recebidas; contadores residuais. \\
T10 & Distância >1 m & 10/\allowbreak 10, 1,12 m; funcional. \\
T11 & Obstáculo passivo & 10/\allowbreak 10 em cada condição; funcional. \\
T12 & Imagem A \ensuremath{\to} B, 64\ensuremath{\times}32 & 256 bytes; hashes iguais; íntegra. \\
T13 & Imagem A \ensuremath{\to} B, 128\ensuremath{\times}64 & 1024 bytes; hashes iguais; íntegra. \\
T14 & Imagem B \ensuremath{\to} A, 64\ensuremath{\times}32 & 256 bytes; hashes iguais; íntegra. \\
\bottomrule\end{tabularx}\end{table}
{\small Critérios transversais: conteúdo correto, número de sequência e END confirmado; em imagens, tamanho, CRC32 e SHA-256 dos pixels. A execução dos 14 ensaios não equivale à aprovação formal de todos os critérios. NR = não registrado.\par}
\setcounter{subsection}{0}\subsection{T01 a T05 — texto e detecção de erros}
\subsubsection*{T01 — Autoteste local}
Executar /teste em A e B. Os registros anteriores indicam CRC=OK, detecção de alteração=OK e 4B6B=OK nos dois nós. Este teste é local e não mede alcance RF.

\subsubsection*{T02 — Texto de A para B}
Zerar contadores, enviar TCD TESTE A PARA B em A e consultar /status nos dois circuitos. O relatório anterior registra texto completo em B e END confirmado em A. A OLED deve apresentar o texto; a fotografia específica não consta da consolidação atual.

\subsubsection*{T03 — Texto de B para A}
Zerar contadores, enviar TCD RESPOSTA B PARA A em B e consultar /status. Foram registrados 5 TX, 3 ACKs e 2 reenvios em B; A contabilizou uma mensagem, sem CRC errado ou duplicatas.

\subsubsection*{T04 — Fragmentação}
Enviar TESTE DE TEXTO LONGO - PACOTES DE 16 BYTES - TCD 2026 a partir de A. A mensagem de 53 bytes foi remontada em B, em quatro DATA (16+16+16+5) e END. A registrou 23 TX, 5 ACKs e 18 reenvios; B registrou 18 duplicatas e uma mensagem. A entrega ocorreu com alta repetição de quadros.

\subsubsection*{T05 — Corrupção e recuperação}
No protocolo, armar /erro em A antes do texto TESTE CRC COM RECUPERACAO. O registro anterior utilizou a mensagem TCD TESTE DE ERRO. O resultado consolidado indica 4 TX, 3 ACKs e 1 reenvio em A; B registrou um CRC errado, nenhuma duplicata e uma mensagem. O descarte sem ACK desencadeou retransmissão.

{\small Situação T01–T05: sucesso relatado, com capturas anteriores reunidas no Anexo A. Distâncias e tempos não registrados.\par}
\setcounter{subsection}{1}\subsection{T06 a T09 — recuperação e estabilidade}
\subsubsection*{T06 — Três erros simulados}
O roteiro prevê /ruido 3 seguido de /lote 2 em A; espera duas mensagens completas. A registrou 10 TX, 4 ACKs e 6 reenvios. B registrou 3 CRCs incorretos, 3 duplicatas e uma mensagem contabilizada. O critério das duas mensagens não está comprovado: evidência parcial, apesar da recuperação registrada.

\subsubsection*{T07 — Perda do ACK de DATA}
Armar /perderack em B e transmitir texto de A (TESTE ACK DATA no roteiro). A: 8 TX, 3 ACKs, 5 reenvios; B: 5 duplicatas, zero CRC errado, uma mensagem. A omissão e a repetição foram registradas; o receptor não acrescentou conteúdo duplicado.

\subsubsection*{T08 — Perda do ACK de END}
Armar /perderend em B, enviar texto de A e observar confirmação de END. O END foi confirmado na terceira tentativa. A: 15 TX, 3 ACKs, 8 reenvios; B: 8 duplicatas, zero CRC errado, uma mensagem. Os contadores incluem atividade residual; não atribuir todos os reenvios à omissão do END.

\subsubsection*{T09 — Lote em curta distância}
Executar /lote 10 em A e consultar /status nos dois nós. B contabilizou 10 mensagens completas, 137 quadros válidos e 107 duplicatas; 30 quadros novos são coerentes com 20 DATA e 10 END. A registrou 145 TX, 32 ACKs, 113 reenvios e 11 mensagens contabilizadas, evidenciando atividade residual. Não calcular taxa de reenvio exclusiva do lote.

{\small Distância e duração dos T06–T09: NR. T06 possui evidência parcial; T07–T09 têm sucesso relatado, com ressalva de contadores no T08/T09.\par}
\setcounter{subsection}{2}\subsection{T10 — comunicação acima de 1 metro}
O procedimento previa posicionar as antenas a 1,10 m, zerar contadores e executar /lote 10 em A. A distância real registrada foi 1,12 m. A confirmou os dez END e B contabilizou dez mensagens. O tempo foi estimado entre 120 e 180 s, não cronometrado.

\begin{table}[H]\centering\small\caption{Contadores do ensaio T10}\label{tab:15}
\rowcolors{2}{tcdlight}{white}\begin{tabularx}{\linewidth}{>{\hsize=1.6392\hsize\linewidth=\hsize\raggedright\arraybackslash}X>{\hsize=0.6804\hsize\linewidth=\hsize\raggedright\arraybackslash}X>{\hsize=0.6804\hsize\linewidth=\hsize\raggedright\arraybackslash}X}
\rowcolor{tcdtablehead} \textcolor{black}{\textbf{Contador}} & \textcolor{black}{\textbf{A (TX)}} & \textcolor{black}{\textbf{B (RX)}} \\
TX /\allowbreak  RX válidos & 49 /\allowbreak  30 & 98 /\allowbreak  49 \\
ACK enviados /\allowbreak  recebidos & 0 /\allowbreak  30 & 98 /\allowbreak  0 \\
Retransmissões /\allowbreak  duplicatas & 19 /\allowbreak  0 & 0 /\allowbreak  19 \\
CRC errado /\allowbreak  4B6B inválido & 0 /\allowbreak  28 & 0 /\allowbreak  0 \\
Mensagens enviadas /\allowbreak  recebidas & 10 /\allowbreak  0 & 0 /\allowbreak  10 \\
\bottomrule\end{tabularx}\end{table}


O lote contém 30 quadros novos (20 DATA + 10 END). Os 49 TX de A correspondem a 30 quadros novos e 19 reenvios; B registrou 49 válidos e 19 duplicatas, além de 98 ACKs enviados (duas cópias por quadro). \textbf{Resultado: sucesso funcional} a 1,12 m; os dez textos integrais não estão visíveis nos registros. Evidências seriais no Anexo A, com o documento original referenciado no Anexo B.

\setcounter{subsection}{3}\subsection{T11 — comunicação com obstáculo físico}
Comparar o enlace na mesma distância informada de 1,12 m, com e sem obstáculo passivo. O roteiro previa condição livre primeiro; na execução, a mochila foi interposta primeiro e a condição livre foi repetida depois. Em ambas, A executou /lote 10 e foram coletados os contadores.

\begin{table}[H]\centering\small\caption{Comparação das condições do ensaio T11}\label{tab:16}
\rowcolors{2}{tcdlight}{white}\begin{tabularx}{\linewidth}{>{\hsize=1.5155\hsize\linewidth=\hsize\raggedright\arraybackslash}X>{\hsize=0.7423\hsize\linewidth=\hsize\raggedright\arraybackslash}X>{\hsize=0.7423\hsize\linewidth=\hsize\raggedright\arraybackslash}X}
\rowcolor{tcdtablehead} \textcolor{black}{\textbf{Indicador}} & \textcolor{black}{\textbf{Sem mochila}} & \textcolor{black}{\textbf{Com mochila}} \\
TX A /\allowbreak  RX válidos B & 44 /\allowbreak  44 & 43 /\allowbreak  43 \\
ACK recebidos A /\allowbreak  enviados B & 30 /\allowbreak  88 & 30 /\allowbreak  86 \\
Reenvios A /\allowbreak  duplicatas B & 14 /\allowbreak  14 & 13 /\allowbreak  13 \\
CRC errado A /\allowbreak  B & 0 /\allowbreak  0 & 0 /\allowbreak  0 \\
4B6B inválido A /\allowbreak  B & 19 /\allowbreak  0 & 22 /\allowbreak  0 \\
Mensagens recebidas /\allowbreak  END & 10 /\allowbreak  10 & 10 /\allowbreak  10 \\
Tempo cronometrado (s) & 200 & 176 \\
\bottomrule\end{tabularx}\end{table}

A conclusão dos dois lotes confirma funcionamento nas duas condições observadas. A diferença de um reenvio e 24 segundos não demonstra que o obstáculo melhorou o enlace: as condições não foram repetidas em uma série controlada. A mochila é um obstáculo passivo, não uma fonte de interferência RF controlada. \textbf{Resultado: sucesso funcional}; os dez textos integrais de cada condição não constam das capturas. Evidências seriais no Anexo A, com o documento original referenciado no Anexo B.

\setcounter{subsection}{4}\subsection{T12 — imagem 64\ensuremath{\times}32 de A para B}
Abrir primeiro o receptor Python B na COM8; iniciar em A, COM20, o envio da imagem4.jpeg em 64\ensuremath{\times}32. O teste foi realizado a 1,12 m, com 112 s de duração. A confirmou END na segunda tentativa e B reconstruiu a imagem.

\par\noindent\begin{minipage}{\linewidth}
\begin{lstlisting}
[B] python receber_imagem_TCD.py COM8
[A] python enviar_imagem_TCD.py COM20 "imagem4.jpeg" --size 64x32
\end{lstlisting}
\end{minipage}\par
\begin{table}[H]\centering\small\caption{Resultados da transmissão de imagem T12}\label{tab:17}
\rowcolors{2}{tcdlight}{white}\begin{tabularx}{\linewidth}{>{\hsize=0.6186\hsize\linewidth=\hsize\raggedright\arraybackslash}X>{\hsize=1.3814\hsize\linewidth=\hsize\raggedright\arraybackslash}X}
\rowcolor{tcdtablehead} \textcolor{black}{\textbf{Parâmetro}} & \textcolor{black}{\textbf{Resultado registrado}} \\
Distância /\allowbreak  duração & 1,12 m /\allowbreak  112 s \\
Bytes recebidos /\allowbreak  esperados & 256 /\allowbreak  256 \\
CRC32 (A/\allowbreak B e PBM) & DAF56504 \\
SHA-256 dos pixels & \texttt{80c42a27\allowbreak 9aaed29d\allowbreak aa665e3a\allowbreak c66efb65\allowbreak 0d83e8d8\allowbreak bb860425\allowbreak 19505986\allowbreak 162c22f7} \\
Arquivo & PBM verificado; PNG registrado no terminal. \\
END & Confirmado na segunda tentativa. \\
\bottomrule\end{tabularx}\end{table}

Após a validação e o salvamento, o receptor registrou erro ClearCommError. Como o erro ocorreu depois da gravação, não invalida o PBM verificado. O hash refere-se aos bytes de pixels, sem o cabeçalho PBM. Não foi registrado o total de retransmissões. \textbf{Resultado: integridade confirmada.} Os terminais e a fotografia estão reunidos no Anexo A e preservados no documento complementar indicado no Anexo B.

\setcounter{subsection}{5}\subsection{T13 — imagem 128\ensuremath{\times}64 de A para B}
Repetir a transferência com a resolução máxima de 128\ensuremath{\times}64. A execução empregou COM20 em A e COM8 em B. Ocorreram muitos erros e a distância foi reduzida de 1,12 m para 0,60 m durante a sessão. O tempo total registrado foi 720 s; END foi confirmado na primeira tentativa.

\par\noindent\begin{minipage}{\linewidth}
\begin{lstlisting}
[B] python receber_imagem_TCD.py COM8
[A] python enviar_imagem_TCD.py COM20 "imagem4.jpeg" --size 128x64
\end{lstlisting}
\end{minipage}\par
\begin{table}[H]\centering\small\caption{Resultados da transmissão de imagem T13}\label{tab:18}
\rowcolors{2}{tcdlight}{white}\begin{tabularx}{\linewidth}{>{\hsize=0.6186\hsize\linewidth=\hsize\raggedright\arraybackslash}X>{\hsize=1.3814\hsize\linewidth=\hsize\raggedright\arraybackslash}X}
\rowcolor{tcdtablehead} \textcolor{black}{\textbf{Parâmetro}} & \textcolor{black}{\textbf{Resultado registrado}} \\
Distância /\allowbreak  duração & 1,12 m \ensuremath{\to} 0,60 m /\allowbreak  720 s \\
Bytes recebidos /\allowbreak  esperados & 1024 /\allowbreak  1024 \\
CRC32 (A/\allowbreak B e PBM) & 2749D731 \\
SHA-256 dos pixels & \texttt{4a9fee35\allowbreak ec9e4c6a\allowbreak 73b0491b\allowbreak c91053fb\allowbreak 49b46081\allowbreak fc9a3218\allowbreak 0939efd8\allowbreak ee7cae3e} \\
Arquivo & PBM verificado; PNG registrado no terminal. \\
END & Confirmado na primeira tentativa. \\
\bottomrule\end{tabularx}\end{table}

A transferência concluiu os 1024 bytes, com CRC32 e SHA-256 coincidentes nos registros e PBM. Não se pode afirmar que toda a imagem percorreu 1,12 m. O total de reenvios não foi registrado. \textbf{Resultado: integridade confirmada.} Os terminais constam do Anexo A e do documento complementar indicado no Anexo B.

\setcounter{subsection}{6}\subsection{T14 — imagem 64\ensuremath{\times}32 de B para A}
Inverter os papéis: A recebeu pela COM20 e B enviou pela COM5, utilizando o comando py. A primeira tentativa foi interrompida na 22ª tentativa do cabeçalho; houve também problema de acesso à COM20. Após diagnóstico por texto e nova sessão, o cabeçalho foi confirmado na quarta tentativa e END na segunda.

\par\noindent\begin{minipage}{\linewidth}
\begin{lstlisting}
[A] python receber_imagem_TCD.py COM20
[B] py enviar_imagem_TCD.py COM5 "imagem4.jpeg" --size 64x32
\end{lstlisting}
\end{minipage}\par
\begin{table}[H]\centering\small\caption{Resultados da transmissão de imagem T14}\label{tab:19}
\rowcolors{2}{tcdlight}{white}\begin{tabularx}{\linewidth}{>{\hsize=0.6186\hsize\linewidth=\hsize\raggedright\arraybackslash}X>{\hsize=1.3814\hsize\linewidth=\hsize\raggedright\arraybackslash}X}
\rowcolor{tcdtablehead} \textcolor{black}{\textbf{Parâmetro}} & \textcolor{black}{\textbf{Resultado registrado}} \\
Distância /\allowbreak  duração & 1,15 m /\allowbreak  120–180 s (estimado). \\
Bytes recebidos /\allowbreak  esperados & 256 /\allowbreak  256 \\
CRC32 (A/\allowbreak B e PBM) & 607BACD6 \\
SHA-256 dos pixels & \texttt{44b987a6\allowbreak 5d851199\allowbreak 8c873caf\allowbreak 6dbfe238\allowbreak a7006d68\allowbreak e8ec7152\allowbreak df4d82ff\allowbreak c2c71a92} \\
Arquivos & PBM verificado; PNG original com os mesmos pixels. \\
END & Confirmado na segunda tentativa. \\
\bottomrule\end{tabularx}\end{table}

O tempo estimado refere-se à tentativa bem-sucedida e não inclui o diagnóstico anterior. A não possui display, portanto a imagem foi reconstruída no computador A. O hash é diferente do T12, mas coincide entre envio e recepção do próprio T14; a causa da diferença entre as conversões não foi determinada. \textbf{Resultado: integridade confirmada.} As evidências estão reunidas no Anexo A.

